% SOURCE docs/release/2026-10-29/source/docs/writing-repair/2026-10-29/src/45358a4637eac946.tex CHAR 0
\chapter*{What the bank should buy first}
\addcontentsline{toc}{chapter}{What the bank should buy first}

A bank considering a better deposit rate must pay for money it already holds
as well as money the offer attracts. A collections team must account for how
today's contact changes tomorrow's information. A discovery algorithm may
return several explanations of the same observations, with opposite implications
for a pricing decision. An apparently clean customer table may still conceal
misrecorded treatment or incorrectly linked accounts. These are four different
reasons why a persuasive numerical answer can fail to improve a decision.

An increase in one account's recorded balance can be a transfer from another
account at the same bank. The funding outcome must aggregate the intended
household and distinguish new external money from internal movement. Competitor
repricing can also change balances during the same period. A lender's threshold
comparison concerns the applicants affected at that margin, while a treasury
shock analysis must separate the external shock from the bank's own response.
These economic distinctions determine which comparison is useful before a
statistical model is chosen.

The proposal is to make a small class of defensible studies easier to conduct
and compare. The first investment should be a comparative integration using
established methods, a precise study description and a report that preserves
uncertainty. A capable economist using existing tools is the reference
alternative. The proposed service earns a place only if the completed study
is better or cheaper at equivalent quality. Broad replacement of the bank's
statistical estate is not justified by the evidence presented here.

That conclusion follows from the examples, rather than from an assumed
advantage of a common interface. In the constructed offer study, a population
balance effect is carried through sampling uncertainty into annual contribution.
The servicing calculation shows why a correctly fitted regression can still
answer the wrong policy question. The graph example gives identical observed
laws with different intervention values. The observation example shows how
additional measurement knowledge can reverse a financial conclusion. Each
case tells us what a service must preserve and what it must decline to infer.

The evidence has different roles. Established causal results justify inference
under stated assumptions. The project's exact derivations and generated studies
expose target errors, weak support and unresolved models. Public data reveal
realistic timing and measurement difficulties. Local interface executions
establish that particular requests retain their meaning through particular
clients. None of these substitutes for an independently reviewed study on
authentic bank information. Their combined value is to identify a useful and
bounded next experiment before undertaking that more expensive work.


